\teachcover{Cyclic groups\\and counting matrices}{1}{
Multiplication in a finite field has a generator
\[
K^\times=\{1,g,g^2,\ldots,g^{q-2}\},\qquad |K|=q.
\]
}{
This lecture proves cyclicity of a finite field's multiplicative group
and then counts ordered bases and determinant fibers. Recall Lagrange's
theorem and the definition of element order. We prove the polynomial
root bound before using it. The group-size-one case is included.
The optional slides compare finite fields with units of prime-power
residue rings; they are not part of the 50-minute main route.
}

\teachslide{The nonzero elements form a group}{2}{
Let $K$ be a field with $q$ elements.
\[
K^\times=K\setminus\{0\},\qquad |K^\times|=N=q-1.
\]
Multiplication is commutative; the identity is $1$;
every element has an inverse.
\[
\key{\ord(a)\mid N,\qquad a^N=1\quad(a\in K^\times).}
\]
}{
Closure follows from absence of zero divisors: the product of two
nonzero elements cannot be zero. Inverses are nonzero. Lagrange
gives the displayed order divisibility. Define order as the least
positive exponent returning to $1$. The task is to find one element
of order $N$, which is stronger than all orders dividing $N$.
}

\teachslide{What cyclic means}{2}{
In $\F_5$, the powers of $2$ are
\[
\begin{array}{c|rrrr}
t&0&1&2&3\\\hline
2^t&1&2&4&3
\end{array}
\qquad 2^4=1.
\]
Every nonzero element is a power of $2$.
\spacer
A group is \key{cyclic} if one element generates the whole group.
}{
Calculate the last power by reducing $8$ modulo $5$. A generator of
a finite group of size $N$ has order $N$, and conversely an element
of order $N$ supplies all $N$ elements by its distinct powers.
There may be several generators. Our aim is to prove at least one
exists in every finite field, not to identify a canonical generator.
}

\teachslide{A polynomial cannot have too many roots}{2}{
Over a field, if $f(a)=0$, then
\[
f(X)=(X-a)h(X).
\]
Any other root $b\ne a$ is a root of $h$:
\[
0=(b-a)h(b)\quad\Longrightarrow\quad h(b)=0.
\]
Induction gives:
\[
\key{\deg f=d,\ f\ne0\quad\Longrightarrow\quad\text{at most }d\text{ roots}.}
\]
}{
Recall polynomial division by $X-a$. The constant nonzero case has
no roots. Each root factors off one linear term; cancellation works
since $b-a$ is a nonzero field element. This proof is the specific
place field arithmetic enters the cyclicity theorem. The root bound
can fail over a ring with zero divisors.
}

\teachslide{Build the order one prime at a time}{2}{
Factor $N=q-1$:
\[
N=\ell_1^{r_1}\cdots\ell_t^{r_t}.
\]
We will construct, for each prime factor $\ell$,
\[
\key{g_\ell\text{ with order }\ell^r,\qquad \ell^r\parallel N.}
\]
Then multiply these elements.

If $N=1$, the group $\{1\}$ is already cyclic.
}{
Explain the double-bar notation: $\ell^r$ divides $N$, but the next
power does not. Distinct primes have relatively prime powers. The
product step requires both coprime orders and commutativity, proved
below. For $K=\F_2$, $N=1$ and $1$ is a generator of the trivial group.
}

\teachslide{Find an element with the required prime factor}{3}{
Fix $\ell^r\parallel N$.

The polynomial
\[
X^{N/\ell}-1
\]
has at most $N/\ell<N$ roots.

It cannot vanish at every element of $K^\times$.
\[
\key{\text{Choose }a\in K^\times\text{ with }a^{N/\ell}\ne1.}
\]
}{
Do not assume a generator to choose $a$: the root bound proves that
at least one of the $N$ nonzero elements fails this equation. The
chosen $a$ can depend on $\ell$, which is fine. Its order must include
the full $\ell$-power of $N$, but the next slide proves the needed
statement directly through powers rather than invoking an order formula.
}

\teachslide{Isolate the prime-power order}{3}{
Set $M=N/\ell^r$ and $g_\ell=a^M$.
\[
g_\ell^{\ell^r}=a^N=1,
\qquad g_\ell^{\ell^{r-1}}=a^{N/\ell}\ne1.
\]
Its order divides $\ell^r$, but does not divide $\ell^{r-1}$.

Every proper divisor of $\ell^r$ divides $\ell^{r-1}$.
\[
\key{\ord(g_\ell)=\ell^r.}
\]
}{
Read both exponent products aloud. If $x^m=1$, its order divides
$m$: divide $m$ by the least positive return time and use minimality
of the remainder. This fact is more precise than Lagrange and also
applies here to a known exponent. The prime-power divisor observation
is what forces the exact order without computing it by trial.
}

\teachslide{Combine two coprime orders}{4}{
Suppose $a,b$ commute, with orders $r,s$ and $\gcd(r,s)=1$.
\[
(ab)^{rs}=1.
\]
If $(ab)^t=1$, put $z=a^t=b^{-t}$. Then
\[
z^r=z^s=1\quad\Longrightarrow\quad z=1
\]
by Bezout's identity. Thus $r\mid t$ and $s\mid t$, so $rs\mid t$.
\[
\key{\ord(ab)=rs.}
\]
}{
This lemma deserves the full four minutes. If $ur+vs=1$, then
$z=z^{ur+vs}=(z^r)^u(z^s)^v=1$; negative exponents use inverses.
Hence $a^t=b^t=1$. The order-divides-return-time fact forces both
divisibilities, whose coprimality gives $rs\mid t$. We already know
$rs$ is a return time, so it is the least. Commutativity is needed
for $(ab)^t=a^tb^t$. The result does not claim that arbitrary
products have the lcm of their factors' orders.
}

\teachslide{One element of full order}{3}{
The prime-power orders $\ell_i^{r_i}$ are pairwise coprime.

By repeatedly applying the lemma,
\[
g=g_{\ell_1}\cdots g_{\ell_t}
\quad\text{has order}\quad
\ell_1^{r_1}\cdots\ell_t^{r_t}=N.
\]
Its $N$ distinct powers exhaust $K^\times$.
\[
\key{\text{The multiplicative group of every finite field is cyclic.}}
\]
}{
Check that a partial product's order remains coprime to the next
prime-power order, so induction applies. An element of order $N$
has exactly $N$ distinct powers from exponent $0$ through $N-1$;
there are only $N$ elements in the group. The proof has used Lagrange,
the polynomial root bound, Bezout, and commutativity, with no general
classification theorem for finite abelian groups.
}

\teachslide{Primitive elements}{2}{
A generator $g$ of $K^\times$ is called a \key{primitive element}.
\[
K^\times=\{1,g,\ldots,g^{q-2}\},\qquad \ord(g)=q-1.
\]
If $K$ has characteristic $p$, then
\[
\key{K=\F_p(g):\text{ the field is generated by one element}.}
\]
Every nonzero element is a power of $g$.
}{
In this finite-field context primitive means generator of the
multiplicative group. Such an element also generates the extension
field, since a field containing it contains all its powers and zero.
The converse need not hold: generating a field extension does not
necessarily mean having multiplicative order $q-1$. In particular,
the Fibonacci root modulo $7$ will have order $16$, not $48$.
}

\teachslide{Your turn: combine orders in $\F_7$}{3}{
In $\F_7^\times$,
\[
\ord(6)=2,\qquad\ord(4)=3.
\]
Verify these orders, then answer:
\begin{enumerate}
\item What is the order of $6\cdot4$?
\item Which residue is therefore a primitive element?
\end{enumerate}
}{
Give a minute of calculation. Students should find $6^2=1$ with
$6\ne1$, and $4^2=2$, $4^3=1$ with $4\ne1$. Coprimality lets us
combine the orders. This is a concrete version of the theorem's
final construction, not a search through all six possible generators.
}

\teachslide{A generator modulo seven}{1}{
\[
6\cdot4=3\pmod7,
\qquad \ord(3)=2\cdot3=6.
\]
\[
\begin{array}{c|rrrrrr}
t&0&1&2&3&4&5\\\hline
3^t&1&3&2&6&4&5
\end{array}
\]
\key{One element generates all six nonzero residues.}
}{
The table verifies the conclusion independently. We used the lemma
to know the order before listing all powers. This ends the cyclicity
portion at minute 28. The rest of this lecture counts matrix groups.
}
